% !TeX root = ../main.tex

\chapter{结论与展望}
\label{chap:conclusion}

\section{主要结论}

本文以自旋群对称性为理论框架，系统研究了交变磁体的分类及其非线性输运性质，并通过 CrSb 单晶的实验测量进行了验证。主要结论如下：

\begin{enumerate}
  \item \textbf{对称性分类方面：} 基于自旋空间群（SSG），给出了交变磁性的严格对称性判据，并依据自旋劈裂的节点结构将交变磁体划分为 $d$ 波、$g$ 波和 $i$ 波三类，揭示了其区别于传统铁磁与反铁磁的本质特征。
  \item \textbf{非线性输运理论方面：} 从量子几何出发，系统推导了 Berry 曲率多极子诱导的非线性霍尔电导率表达式，并结合磁点群的对称性分析，给出了 122 个三维磁点群中主导反常霍尔响应阶次的完整列表。
  \item \textbf{CrSb 的对称性约束方面：} 针对磁点群 $6'/m'mm'$，严格推导了各阶电导率张量的对称性允许形式：一阶无本征反常霍尔效应，二阶完全消失，三阶横向非线性霍尔响应仅在 $x$ 方向驱动、$z$ 方向测量时非零。
  \item \textbf{实验验证方面：} 旋转角度依赖测量排除了几何泄露与一阶本征霍尔信号；三阶横向电压仅出现在 $XZ$ 器件配置中，与对称性预测一致；多波混频实验给出了约 20 条清晰的混频谱线，验证了 BCQ 四极矩驱动的瞬时（$>100\,\mathrm{kHz}$）非线性响应特性。
\end{enumerate}

\section{展望}

尽管本文建立了交变磁体非线性输运的对称性分析框架，并在 CrSb 中完成了初步的实验验证，仍有若干方向值得进一步探索：

\begin{itemize}
  \item 本文对称性分析目前限于共线交变磁体，非共线、非共面的交变磁序有望带来更丰富的非线性响应，值得系统推广。
  \item 在材料方面，目前实验仅聚焦于 CrSb。将对称性分析方法推广至其他 $d$ 波、$g$ 波和 $i$ 波交变磁体候选材料，有助于全面检验理论框架。
  \item 交变磁体兼具零净磁矩与时间反演破缺的能带特征，在超导-磁近邻效应、反铁磁自旋电子学器件以及 THz 频段信号转换等方面具备独特的应用潜力，值得进一步挖掘。
\end{itemize}
